\chapter{Run 3: one network for every standoff}\label{ch:run3}

Chapters~\ref{ch:floor} to~\ref{ch:magic_results} end on 8 November, with run 3 still training. This chapter and the
eight after it cover 8 to 16 November. They are written the same way as the rest: every idea is rebuilt from the
start, and every grey block is copied from the repository.

\section{What run 3 changes}

Run 2 was trained on standoffs of $3\pm0.6$\um. You asked whether one network could cover \emph{any} standoff. The
network never sees the height, so the architecture needs no change. Only one line of the data generator differs: how
the standoff of each training crop is drawn.

\begin{code*}{magic_data.py}
        if cfg.get('z_uniform'):                  # run 3: any standoff in [lo, hi]
            zs.append(float(rng.uniform(*cfg['z_uniform'])))
            continue
        z = rng.normal(cfg['z_mean'], cfg['z_std'])
        if abs(z - cfg['z_mean']) <= cfg['z_trunc'] * cfg['z_std']:
            zs.append(float(z))
\end{code*}

Run 2 takes the bottom three lines: a normal distribution around 3\um, cut at three standard deviations. Run 3 takes the
top branch with \code{--z-uniform 1 10}: every height from 1 to 10\um\ is equally likely. Everything else is identical.

\section{The result}

On the simulated validation layouts, with the target always at 0.5\um:

\begin{center}
\small
\begin{tabular}{lcccccccccc}
\toprule
Input standoff (\textmu m) & 1 & 2 & 3 & 4 & 5 & 6 & 7 & 8 & 10 \\
\midrule
Run 2 ($3\pm0.6$\um) & 22.8\% & 14.2\% & 15.1\% & 15.4\% & 25.4\% & 48.1\% & 67.9\% & 80.5\% & 92.0\% \\
Run 3 (1--10\um) & 35.3\% & 26.1\% & 22.7\% & 22.1\% & 26.1\% & 31.2\% & 33.7\% & 36.1\% & 56.8\% \\
\bottomrule
\end{tabular}
\end{center}

And on COMSOL at 3\um\ to 0.5\um:

\begin{center}
\begin{tabular}{lccc}
\toprule
 & 15 pairs & C1-family coils & \code{IC_4um} \\
\midrule
Run 2 & 27.0\% & 26.3\% & 29.6\% \\
Run 3 & 34.8\% & 30.4\% & 52.0\% \\
\bottomrule
\end{tabular}
\end{center}

\begin{keyidea}[What run 3 shows]
\begin{itemize}
\item \textbf{A wide range has a price.} Inside run 2's range (2--4\um) run 3 is 7 to 12 points worse. A network that must
be ready for any blur is less sure about each one.
\item \textbf{It buys reach.} From 5 to 8\um\ run 3 holds at 26--36\% where run 2 collapses.
\item \textbf{It breaks at the edge of its own range.} At 10\um\ it jumps to 56.8\%. Remember this number: it is the reason
the 10\um\ runs of Chapter~\ref{ch:ten} are trained with 10\um\ in the \emph{middle} of their range.
\end{itemize}
\end{keyidea}

\section{Why the IC coil fails}

On \code{IC_4um} run 3 scores 52\%, and by component it is 31 / 69 / 56\% for $B_x$ / $B_y$ / $B_z$. That is not noise. It has
a cause, and the cause is the ambiguity of Chapter~\ref{ch:question}.

The IC coil has pairs of wires close together that carry current in \emph{opposite} directions. Far above such a pair
the two fields nearly cancel. What survives is weak and smooth. But ``weak and smooth'' is also exactly what a single
wire looks like from far away. So a tile that sees such a pair has two explanations:
\begin{itemize}
\item two opposite wires, close to the sensor; or
\item one source, far from the sensor, which must then be made much stronger and sharper to explain the same field.
\end{itemize}
Run 2 only knows heights near 3\um, so it has no choice but the first. Run 3 was told that anything from 1 to 10\um\ is
possible, and on these tiles it picks ``far''. It then over-amplifies, and the 0.5\um\ field comes out wrong.

There is a second ingredient. MAGIC-UNet works tile by tile, and each $12.8$\um\ tile makes this decision \emph{alone}.
Nothing tells neighbouring tiles that the whole image was measured at one height.

\begin{careful}[What is established and what is not]
The pattern (the failure sits on the closely spaced opposite wires, and different tiles imply different heights) was
measured on the server. The explanation is the best one we have, and Chapter~\ref{ch:height} tests one consequence of
it: if the only problem were the per-tile height, then forcing \emph{one} height on the whole image should rescue run 3.
It does not.
\end{careful}

\begin{defend}[If your professor asks: so can one network handle an unknown standoff?]
Inside a narrow band, yes: run 2 works from 2 to 4\um\ without being told the height. Across a wide band it works on
simple wires and fails where the wires are ambiguous, because the field alone cannot separate ``close and weak'' from
``far and strong''. That is the height ambiguity showing up inside a trained network.
\end{defend}
